\section{Introduction}\label{sec:introduction}

\IEEEPARstart{T}{HERE} was a time when serving large language models (LLMs) felt like a compute problem. That has changed. Once parameter counts run past what fits on a single Graphics Processing Unit (GPU), the question stops being how many TFLOPS you can extract and starts being how data moves: between GPU Video RAM (VRAM), host Dynamic RAM (DRAM), and Non-Volatile Memory express (NVMe) storage. End-to-end latency is then the sum of capacity ceilings, cross-tier bandwidth, transfer cost, and synchronization. Raw arithmetic still matters, but mostly as a baseline---it is no longer where the variance lives.

This shift hits industrial informatics first. Inference there is not optional or eventual; it is sitting inside production. Digital twin synchronization. Predictive maintenance over Industrial Internet of Things (IIoT) sensor streams. Real-time inspection on the line. These are closed-loop cyber--physical systems (CPS), and the inference engine has a different job than it does in a research lab. It is not asked to be fast on average. It is asked to be predictably fast, scan-cycle aligned, and inside the energy envelope, every minute the line is running. Industrial workloads are also the kind of thing that misbehave on demand---a sensor goes anomalous and you suddenly have a burst, several models contending for the same accelerator, and an SLA you can't miss. Determinism and long-running stability stop being nice-to-haves and become hard requirements.

Latency in this regime is not just a quality-of-service question. It is a safety question.

Bounded response times have to line up with industrial real-time standards: Time-Sensitive Networking (TSN) windows, the cyclic Programmable Logic Controller (PLC) execution semantics defined by IEC~61131-3. When inference sits inside the control loop, it has to finish inside a reserved sub-window of the scan cycle---otherwise the input--logic--output ordering breaks down. The 2\,ms budget we work with later is for bounded warm-state inference (scoring, short-token continuation), not autoregressive generation of a 17B model. By holding tail latency and swap intensity ($\beta < 1$) within bounds, RAMSES finishes inside the reserved TSN window and the scan deadline. From the IEC~61508 angle, systematic timing faults feed the residual failure rate $\lambda_{\mathrm{res}}$ that characterizes dangerous undetected failures; we therefore define $P_{\mathrm{timing}} = \Pr(T_{\mathrm{total}} > T_{\mathrm{budget}})$ as the timing-related fault probability. Holding $\beta < 1$ together with $\alpha \ge \alpha_{\mathrm{critical}}$ keeps $T_{\mathrm{swap}}$ below $T_{\mathrm{comp}} + T_{\mathrm{mem}}$, which lowers $P_{\mathrm{timing}}$ and trims the timing-induced share of $\lambda_{\mathrm{res}}$.

Coordinating memory more aggressively is, at this point, surprisingly easy to do badly. The same offloading or caching policy can speed one configuration up and slow another one down. Why? Because the system is in different regimes, and a different latency component dominates in each. Capacity-limited execution: the working set does not fit in fast tiers. I/O-limited execution: cross-tier traffic saturates storage bandwidth. Coordination-dominated execution: compute, locality, transfer, and sync sit at comparable scale. Only in the third regime does coordination buy anything. For industrial CPS and IIoT, though, the regime structure of hierarchical memory has not been written down in a way that supports deterministic operation, which is the gap we set out to close.

Most prior work attacks the problem from one direction at a time. Model-level work (quantization, pruning) shrinks the footprint. System-level work (CPU/GPU offloading, caching) shuffles data between tiers. Recent papers in IEEE Transactions on Industrial Informatics study resource allocation, cloud--edge orchestration, and latency-aware scheduling, but they reason at the network layer or about distributed task placement and treat each node's memory as a black box~\cite{tao2019digitaltwin,cao2021edgecps,chen2021drl,wu2021dnninference,lu2021dtfl,chekired2018fogscheduling,deng2020edgerl}. We open the box. The question we focus on is how VRAM, DRAM, and NVMe interact inside one serving node, and we model the cross-tier transfer formally rather than empirically.

The framework is called \textbf{RAMSES} (Regime-Aware Memory System for Efficient Serving). It is a lightweight orchestrator that treats VRAM, DRAM, and NVMe as one hierarchy and enforces regime-aware control. Conventional serving stacks chase average throughput; RAMSES is sized for a different goal---bounded delay and long-running stability, both required by industrial deployments. The capacity and I/O pressure ratios let us write down, in closed form, when hierarchical coordination remains stable across sustained industrial workloads.

To check whether all of this actually holds up, we run RAMSES on language and vision workloads representative of industrial usage, against strong offloading baselines and modern inference engines. Loading time and latency drop. Peak memory use drops. The system stays stable across a 72-hour stress run. Throughput-per-watt improves, as measured by GPU telemetry. Ablations indicate that the three modules pull in concert, not in isolation.

What we contribute, concretely:
\begin{itemize}
\item \emph{Regime-Aware Analytical Model:} A multi-tier latency decomposition with explicit boundary conditions that distinguish capacity-limited, I/O-limited, and coordination-dominated regimes for industrial AI infrastructure.
\item \emph{Energy--Latency Coupling Analysis:} A direct analytical link between cross-tier transfer dynamics and throughput-per-watt under sustained SLA constraints.
\item \emph{Lightweight Multi-Tier Orchestration:} The RAMSES design and implementation, with Non-Uniform Memory Access (NUMA)-aware and swap-aware regime control across VRAM, DRAM, and NVMe.
\item \emph{Industrial-Scale Validation:} Long-running workloads showing better SLA stability, energy efficiency, and robustness than strong baselines.
\end{itemize}
